% GENERATED FILE. Edit canonical lessons, then run npm run generate:books.
% canonical-source-hash: fnv1a64:9ebe4751119ff7ed
% canonical-lessons: KA-C176-saluvagi, KA-C176-illadiddare, KA-C176-adaru, KA-C176-ondu-vele, KA-C176-badalu

\chapter{For the sake of, Otherwise, Even so, In case, Instead}
\label{ch:ka-for-the-sake-of-otherwise-even-so-in-case-instead}
\hlchaptermodality{176}

\begin{chapteropening}
\textbf{By the end of this chapter:} \emph{I can say for the sake of, otherwise, even so, in case and instead.}

Complete the last lesson of chapter 176: I can say for the sake of, otherwise, even so, in case and instead.
\end{chapteropening}

\section[\texorpdfstring{\kn{ಸಲುವಾಗಿ} (saluvāgi)}{ಸಲುವಾಗಿ (saluvāgi)}]{\kn{ಸಲುವಾಗಿ} (saluvāgi) — for the sake of, for}

\label{lesson:KA-C176-saluvagi}



\begin{quote}
Before the new one: say the Kannada for on the other side, then the Kannada for on this side.
\end{quote}

\subsection*{You'll want to know: \kn{ಸಲುವಾಗಿ}}
\textbf{\kn{ಸಲುವಾಗಿ}} — \emph{saluvāgi} — \textquotedblleft{}for the sake of, for\textquotedblright{}.

\begin{grammarlens}[title={What you've built}]
One more word for this chapter.
\end{grammarlens}

\subsection*{Your turn}
\begin{itemize}
\raggedright
  \item \emph{Say it:} \emph{saluvāgi}
  \item \emph{Say it:} \emph{saluvāgi}, once more
  \item \emph{Say it:} say \emph{saluvāgi}
  \item \emph{Recall:} say the Kannada for on the other side, then the Kannada for on this side, then say \emph{saluvāgi} again
\end{itemize}

\begin{tcolorbox}[breakable,colback=teal!4,colframe=teal!35!black,title={Before you move on}]
What does \kn{ಸಲುವಾಗಿ} mean? (\textquotedblleft{}For the sake of, for\textquotedblright{}.) Say it once more.
\end{tcolorbox}
\section[\texorpdfstring{\kn{ಇಲ್ಲದಿದ್ದರೆ} (illadiddare)}{ಇಲ್ಲದಿದ್ದರೆ (illadiddare)}]{\kn{ಇಲ್ಲದಿದ್ದರೆ} (illadiddare) — otherwise}

\label{lesson:KA-C176-illadiddare}



\begin{quote}
Before the new one: say the Kannada for on this side, then the Kannada for for the sake of.
\end{quote}

\subsection*{You'll want to know: \kn{ಇಲ್ಲದಿದ್ದರೆ}}
\textbf{\kn{ಇಲ್ಲದಿದ್ದರೆ}} — \emph{illadiddare} — \textquotedblleft{}otherwise\textquotedblright{}.

\begin{grammarlens}[title={What you've built}]
One more word for this chapter.
\end{grammarlens}

\subsection*{Your turn}
\begin{itemize}
\raggedright
  \item \emph{Say it:} \emph{illadiddare}
  \item \emph{Say it:} \emph{illadiddare}, once more
  \item \emph{Say it:} say \emph{illadiddare}
  \item \emph{Recall:} say the Kannada for on this side, then the Kannada for for the sake of, then say \emph{illadiddare} again
\end{itemize}

\begin{tcolorbox}[breakable,colback=teal!4,colframe=teal!35!black,title={Before you move on}]
What does \kn{ಇಲ್ಲದಿದ್ದರೆ} mean? (\textquotedblleft{}Otherwise\textquotedblright{}.) Say it once more.
\end{tcolorbox}
\section[\texorpdfstring{\kn{ಆದರೂ} (ādarū)}{ಆದರೂ (ādarū)}]{\kn{ಆದರೂ} (ādarū) — even so, nevertheless}

\label{lesson:KA-C176-adaru}



\begin{quote}
Before the new one: say the Kannada for for the sake of, then the Kannada for otherwise.
\end{quote}

\subsection*{You'll want to know: \kn{ಆದರೂ}}
\textbf{\kn{ಆದರೂ}} — \emph{ādarū} — \textquotedblleft{}even so, nevertheless\textquotedblright{}.

\begin{grammarlens}[title={What you've built}]
One more word for this chapter.
\end{grammarlens}

\subsection*{Your turn}
\begin{itemize}
\raggedright
  \item \emph{Say it:} \emph{ādarū}
  \item \emph{Say it:} \emph{ādarū}, once more
  \item \emph{Say it:} say \emph{ādarū}
  \item \emph{Recall:} say the Kannada for for the sake of, then the Kannada for otherwise, then say \emph{ādarū} again
\end{itemize}

\begin{tcolorbox}[breakable,colback=teal!4,colframe=teal!35!black,title={Before you move on}]
What does \kn{ಆದರೂ} mean? (\textquotedblleft{}Even so, nevertheless\textquotedblright{}.) Say it once more.
\end{tcolorbox}
\section[\texorpdfstring{\kn{ಒಂದು} \kn{ವೇಳೆ} (ondu vēḷe)}{ಒಂದು ವೇಳೆ (ondu vēḷe)}]{\kn{ಒಂದು} \kn{ವೇಳೆ} (ondu vēḷe) — in case, if by chance}

\label{lesson:KA-C176-ondu-vele}



\begin{quote}
Before the new one: say the Kannada for otherwise, then the Kannada for even so.
\end{quote}

\subsection*{You'll want to know: \kn{ಒಂದು} \kn{ವೇಳೆ}}
\textbf{\kn{ಒಂದು} \kn{ವೇಳೆ}} — \emph{ondu vēḷe} — \textquotedblleft{}in case, if by chance\textquotedblright{}.

\begin{grammarlens}[title={What you've built}]
One more word for this chapter.
\end{grammarlens}

\subsection*{Your turn}
\begin{itemize}
\raggedright
  \item \emph{Say it:} \emph{ondu vēḷe}
  \item \emph{Say it:} \emph{ondu vēḷe}, once more
  \item \emph{Say it:} say \emph{ondu vēḷe}
  \item \emph{Recall:} say the Kannada for otherwise, then the Kannada for even so, then say \emph{ondu vēḷe} again
\end{itemize}

\begin{tcolorbox}[breakable,colback=teal!4,colframe=teal!35!black,title={Before you move on}]
What does \kn{ಒಂದು} \kn{ವೇಳೆ} mean? (\textquotedblleft{}In case, if by chance\textquotedblright{}.) Say it once more.
\end{tcolorbox}
\section[\texorpdfstring{\kn{ಬದಲು} (badalu)}{ಬದಲು (badalu)}]{\kn{ಬದಲು} (badalu) — instead (of)}

\label{lesson:KA-C176-badalu}



\begin{quote}
Before the new one: say the Kannada for even so, then the Kannada for in case.
\end{quote}

\subsection*{You'll want to know: \kn{ಬದಲು}}
\textbf{\kn{ಬದಲು}} — \emph{badalu} — \textquotedblleft{}instead (of)\textquotedblright{}.

\begin{grammarlens}[title={What you've built}]
One more word for this chapter.
\end{grammarlens}

\subsection*{Your turn}
\begin{itemize}
\raggedright
  \item \emph{Say it:} \emph{badalu}
  \item \emph{Say it:} \emph{badalu}, once more
  \item \emph{Say it:} say \emph{badalu}
  \item \emph{Recall:} say the Kannada for even so, then the Kannada for in case, then say \emph{badalu} again
\end{itemize}

\begin{tcolorbox}[breakable,colback=teal!4,colframe=teal!35!black,title={Before you move on}]
What does \kn{ಬದಲು} mean? (\textquotedblleft{}Instead (of)\textquotedblright{}.) Say it once more.
\end{tcolorbox}
